\section{从 Base Model 到 Chatbot：三种监督信号}

前一章确定了证据边界，本章先建立共同语言。所谓“训练 chatbot”至少包含 demonstration learning、preference learning 与 policy optimization 三层；它们使用不同数据结构，解决不同问题，也需要不同评价方法。

\subsection{SFT、Reward Model 与 RLHF}

\term{SFT（supervised fine-tuning，监督微调）}把 prompt $x$ 与人工/合成 completion $y$ 作为监督对，最小化 next-token cross-entropy：
\[
\mathcal{L}_{\mathrm{SFT}}(\theta)
=-\sum_{t=1}^{|y|}\log \pi_{\theta}(y_t\mid x,y_{<t}).
\]
其中 $\theta$ 是 chatbot 参数，$y_t$ 是第 $t$ 个目标 token，$y_{<t}$ 是已给定前缀。SFT 学的是“像示范那样回答”，并不直接学会比较两个答案谁更好。

\lecturefigure{slide-04.jpg}{InstructGPT 式训练流程：SFT 负责 helpfulness 起点}{官方幻灯片 p.4}

\begin{knowledgebox}{读图：三步的 label 结构不同}
第一步由人写 prompt 与 demonstration；第二步由模型生成多个回答、人只做排序；第三步用 learned reward 更新 policy。先看每步谁产生文本，再看人类劳动从“写完整答案”变成“比较答案”，即可理解为什么三阶段的数据规模、噪声和成本不同。
\end{knowledgebox}

\term{reward model（奖励模型，RM）}把 $(x,y)$ 映射为标量 $r_{\phi}(x,y)$。对 preferred answer $y_w$ 与 rejected answer $y_l$，常用 Bradley--Terry pairwise loss：
\[
\mathcal{L}_{\mathrm{RM}}(\phi)
=-\log \sigma\!\left(r_{\phi}(x,y_w)-r_{\phi}(x,y_l)\right).
\]
其中 $\phi$ 是 reward model 参数，$\sigma$ 是 sigmoid；模型只需让 winner 分数高于 loser，并不要求奖励值具有绝对物理意义。

\lecturefigure{slide-05.jpg}{Preference ranking 与 RLHF 增加 harmlessness/价值约束}{官方幻灯片 p.5}

\begin{importantbox}{RLHF 的目标不是“最大化 reward 就完事”}
\term{RLHF（reinforcement learning from human feedback）}用 learned reward 优化 policy，同时通常用 KL penalty 约束它不要偏离 reference model：
\[
\max_{\theta}\;\mathbb{E}_{y\sim\pi_{\theta}(\cdot\mid x)}
\left[r_{\phi}(x,y)-\beta\log\frac{\pi_{\theta}(y\mid x)}{\pi_{\mathrm{ref}}(y\mid x)}\right].
\]
这里 $\beta$ 控制“追求偏好”与“保持原模型行为”之间的张力。Reward hacking、distribution shift 与 RM 误差都可能让更高分不等于更好答案。
\end{importantbox}

\teachervoice{讲者把 step 1 解释为“让模型先成为 helpful chatbot”，再把 step 2/3 看成加入 guardrails。这个说法是课堂叙事，不意味着 SFT 只影响 helpfulness、RLHF 只影响 harmlessness；现实中三阶段会同时改变风格、事实性与拒答行为。}

\subsection{术语消化：三类样本不要混用}

下面的区分是后文读图的最小前置知识。若把 demonstration、preference pair 与 policy rollout 都叫“训练数据”，就无法解释为什么相同数量的样本成本和信息量完全不同。

\begin{knowledgebox}{三类监督信号}
\begin{tabularx}{\textwidth}{@{}L{0.18\textwidth}L{0.25\textwidth}X@{}}
\toprule
对象 & 样本结构 & 学到什么 \\
\midrule
SFT & prompt + target completion & 模仿目标回答的 token distribution \\
Reward model & prompt + chosen/rejected responses & 学习相对偏好排序 \\
RL policy & prompt + sampled response + scalar reward & 让 policy 更常产生高 reward response \\
\bottomrule
\end{tabularx}
\end{knowledgebox}

\subsection{本章小结}

SFT、reward modeling 与 RLHF 使用三种不同 label。后面讨论 10K demonstrations、20K dialogues 或 100K preference examples 时，必须先问它们属于哪一种样本结构，而不能直接比较数字大小。

